
\documentclass[11pt,a4paper]{article}
\usepackage[a4paper,margin=2.0cm]{geometry}
\usepackage[T1]{fontenc}
\usepackage[utf8]{inputenc}
\usepackage[spanish,es-noquoting]{babel}
\usepackage{lmodern,microtype}
\usepackage{amsmath,amssymb,amsthm,mathtools}
\usepackage{booktabs,longtable,array,enumitem}
\usepackage{xcolor,hyperref}
\pagecolor{white}
\color{black}
\setlength{\parindent}{0pt}
\setlength{\parskip}{0.72em}

\newcommand{\APP}{\mathrm{APP}}
\newcommand{\TPK}{\mathrm{TPK}}
\newcommand{\HMT}{\mathrm{HMT}}
\newcommand{\KR}{K_{R12}^{\rm CAN}}
\newcommand{\Witt}{W_{12}}
\newcommand{\Golay}{\mathcal C_{12}}
\newcommand{\Leech}{\Lambda_{\rm Leech}}
\newcommand{\Monster}{\mathbb M}
\newcommand{\eqledger}{\equiv_{\rm ledger}}

\title{\bfseries Ventana N75 --- Síntesis extensa del cierre holográfico del infinito\\
\large APP, triadas, \(\pi,e,\varphi,\alpha\), cilindros, nueve puertas, Witt/Golay y Moonshine}
\author{Ventana de trabajo HMT}
\date{}

\begin{document}
\maketitle

\begin{abstract}
Esta ventana reúne el hilo técnico y conceptual consolidado: desde la APP \(9\times9\) y la célula activa \(729=3^6\), hasta la lectura triádica de \(\pi,e,\varphi\), el cierre de \(\alpha\) por carry base--\(1000\), la estructura Witt/Golay, el puente Moonshine y la reformulación final de la generación profunda como máquina de nueve puertas con memoria de carry. La tesis es que HMT no genera cifras como una expansión plana, sino eventos de frontera. El número HMT no es sólo una cinta decimal: es una célula APP--729 que se autolee hacia dentro como cilindro ternario/decimal y hacia fuera como simetría excepcional. La raíz vuelve; el borde recuerda.
\end{abstract}

\section{Tesis general}

La formulación madura del programa puede escribirse así:
\[
\boxed{
\text{un número HMT es una célula discreta de borde que se actualiza como lectura real.}
}
\]
Esto desplaza el estatuto del número. El número no aparece primero como punto de la recta real, sino como estructura finita de compatibilidades:
\[
\text{fase},\quad
\text{borde},\quad
\text{carry},\quad
\text{supervivencia},\quad
\text{lectura archimedeana}.
\]
La lectura real es resultado de un cierre. No se niega la completitud de \(\mathbb R\); se usa. Pero HMT propone que la recta real es el lector final de una arquitectura previa de borde:
\[
\boxed{
\text{HMT no sustituye el continuo; lo hace emerger como lectura estabilizada.}
}
\]
Y, en el lenguaje de la serie:
\[
\boxed{
\text{la precisión no añade cifras: estrecha cilindros de supervivencia.}
}
\]
Medir más significa restringir ramas compatibles sin destruir la cola. La cola no es ruido: es memoria de borde.

\section{APP como célula primaria}

La APP baja no es sólo tabla de multiplicar módulo \(9\). Es una célula de fase sobre:
\[
\{1,\ldots,9\}\times\{1,\ldots,9\}.
\]
Tiene dos lectores:
\[
\APP_{\times}(i,j)=\operatorname{dr}_9(ij),
\qquad
\APP_{+}(i,j)=\operatorname{dr}_9(i+j).
\]
Suma y producto no son operaciones separadas: en APP actúan como dos hojas de una misma superficie finita. La suma acumula; el producto pliega. La raíz digital no destruye información sin más: devuelve la operación al borde de la unidad.

La inversión ordinal
\[
\rho(i)=10-i
\]
produce reversibilidad de lectura. Conceptualmente:
\[
\boxed{
\text{multiplicación invertida}
\longrightarrow
\text{defecto aditivo}.
}
\]
Esa es una primera señal de que APP no es aritmética ordinaria, sino geometría de cierre.

\section{De la APP plana a la célula activa \(729\)}

La APP plana tiene \(9\times9=81\) posiciones, pero la lectura de triadas obliga a elevar la célula a:
\[
9^3=729=3^6.
\]
Cada triada decimal \(abc\) vive en:
\[
\{0,\ldots,9\}^3,
\qquad
10^3=1000.
\]
El bulk activo sin ceros es:
\[
9^3=729.
\]
La corona con al menos un cero es:
\[
1000-729=271.
\]
Por tanto:
\[
\boxed{
1000=729+271.
}
\]
No es un adorno decimal. Es la interfaz entre célula triádica activa y corona archimedeana de lectura.

Además:
\[
1000\equiv1\pmod9.
\]
Cada triada decimal vuelve a la unidad de fase módulo \(9\), pero vuelve con carry:
\[
\boxed{
\text{la raíz vuelve; el borde recuerda.}
}
\]

\section{Semillas \(w_6\) de \(\pi,e,\varphi\)}

El nivel semilla consolidado es:
\[
w_6(\pi)=010211,
\]
\[
w_6(e)=201101,
\]
\[
w_6(\varphi)=121200.
\]
Estas semillas no son cifras aisladas. Son regiones iniciales:
\[
\pi=\text{cierre circular},
\qquad
e=\text{propagación/exponencialidad},
\qquad
\varphi=\text{autoescala/proporción}.
\]
La lectura fuerte es:
\[
\boxed{
\pi,e,\varphi
\text{ son regiones de operación distintas en la misma célula de borde.}
}
\]
Desde ellas se obtienen prefijos triádicos:
\[
\pi:\quad 141|592,\quad 141|592|653,\quad 141|592|653|589,
\]
\[
e:\quad 718|281,\quad 718|281|828,\quad 718|281|828|459,
\]
\[
\varphi:\quad 618|033,\quad 618|033|988,\quad 618|033|988|749.
\]
La cautela correcta es:
\[
\boxed{
\text{esto no debe venderse solo como generación infinita completa.}
}
\]
Pero sí como:
\[
\boxed{
\APP/\TPK\text{ abre correctamente las regiones iniciales de }\pi,e,\varphi.
}
\]

\section{Corrector canónico y cierre de \(\alpha\)}

El corrector canónico es:
\[
\boxed{
\KR
=
(234,543,140,729,659,824,621,058,914,794,146,601).
}
\]
La ruta canónica lo fija antes del cálculo de \(\alpha\):
\[
\APP/\TPK/\kappa27/SU\text{-}12
\Longrightarrow
\KR
\Longrightarrow
\text{carry base--}1000
\Longrightarrow
\alpha.
\]
La comparación con CODATA es auditoría externa, no input.

Con las triadas:
\[
\pi:\ 141|592|653|589|793|238|462|643|383|279|502|884,
\]
\[
e:\ 718|281|828|459|045|235|360|287|471|352|662|497,
\]
\[
\varphi:\ 618|033|988|749|894|848|204|586|834|365|638|117,
\]
se obtiene:
\[
\boxed{
\alpha:\ 007|297|352|569|283|800|997|285|105|472|380|663.
}
\]
Conceptualmente:
\[
\boxed{
\alpha
=
\text{contracción de cierre de }(\pi,e,\varphi)\text{ corregida por }\KR.
}
\]
En términos de canal:
\[
\sigma=(1,1,-1),
\]
de modo que \(\alpha\) aparece como lectura firmada:
\[
\pi+e-\varphi-\KR.
\]
Así:
\[
\boxed{
\pi,e,\varphi
\text{ son canales de operación; }\alpha\text{ es cierre de canal.}
}
\]

\section{Witt, Golay y geometría de cierre}

El diseño de Witt aparece como:
\[
\boxed{
\Witt=S(5,6,12).
}
\]
El código ternario extendido asociado tiene:
\[
[12,6,6]_3,
\qquad
|\Golay|=3^6=729,
\]
\[
1+264y^6+440y^9+24y^{12}.
\]
Los soportes de peso \(6\) producen:
\[
264/2=132
\]
hexadas. El automorfismo:
\[
|{\rm Aut}(W_{12})|=95040
\]
identifica:
\[
{\rm Aut}(W_{12})\cong M_{12}.
\]
No es decoración. Es una estructura finita excepcional que rigidifica el borde HMT.

La matriz derecha del generador ternario:
\[
G=[I_6\mid A_W]
\]
satisface:
\[
\boxed{
A_W^2=-I\pmod3,
\qquad
A_W^{-1}=-A_W.
}
\]
Esto convierte \(A_W\) en estructura compleja ternaria.

La preimagen Witt de \(\varphi\) es:
\[
\boxed{
q_\varphi=w_6(\varphi)A_W^{-1}=122120.
}
\]
Y:
\[
q_\varphi A_W=w_6(\varphi).
\]
La autoescala \(\varphi\) no sólo aparece como valor o región visible. También aporta coordenada de borde:
\[
\boxed{
\varphi\longrightarrow q_\varphi.
}
\]
La lectura:
\[
\boxed{
\text{la autoescala aporta coordenadas; el cierre aporta cargas.}
}
\]

\section{Moonshine desde APP}

En APP:
\[
S_+(9)=405,
\qquad
S_{\times}(9)=459,
\qquad
\Delta(9)=54.
\]
También:
\[
744=\Sigma(9)-120,
\]
y:
\[
\boxed{
196884=54(3645+1).
}
\]
Esto conecta APP con la aritmética de Moonshine. La cadena conceptual:
\[
W_{12}\to W_{24}\to\Leech\to V^\natural\to\Monster.
\]
La formulación sobria:
\[
\boxed{
\text{HMT cae en la puerta finita de la escalera Golay--Witt--Leech--VOA--Monster.}
}
\]
No se afirma que HMT ``demuestre el Monstruo''. Se afirma que:
\[
\boxed{
\text{el borde HMT pertenece naturalmente al ambiente excepcional que Moonshine completa.}
}
\]

\section{Del dígito al cilindro}

El error corregido fue tratar un prefijo ternario como punto. El prefijo ternario de longitud \(L\) es:
\[
\left[
\frac{N}{3^L},
\frac{N+1}{3^L}
\right).
\]
El prefijo decimal de \(K\) triadas es:
\[
\left[
\frac{D}{1000^K},
\frac{D+1}{1000^K}
\right).
\]
La compatibilidad correcta es:
\[
\boxed{
\left[
\frac{N}{3^L},
\frac{N+1}{3^L}
\right)
\cap
\left[
\frac{D}{1000^K},
\frac{D+1}{1000^K}
\right)
\neq\varnothing.
}
\]
Equivalente:
\[
\boxed{
N1000^K<(D+1)3^L,
}
\]
\[
\boxed{
(N+1)1000^K>D3^L.
}
\]
Esto formaliza:
\[
\boxed{
0_{\rm borde}\neq9_{\rm activo}.
}
\]
El borde no se decide por el extremo inferior, sino por el intervalo completo.

\section{Estado de cilindro y transición exacta}

Definimos:
\[
A=N1000^K-D3^L,
\qquad
B=(N+1)1000^K-D3^L.
\]
La supervivencia:
\[
\boxed{
A<3^L,\qquad B>0.
}
\]

Al añadir un bloque ternario:
\[
u\in\{0,\ldots,728\},
\]
se tiene:
\[
N'=729N+u.
\]
Si no entra nueva triada decimal:
\[
\Delta K=0,
\]
entonces:
\[
\boxed{
A'=729A+u1000^K.
}
\]
Si entra nueva triada:
\[
\Delta K=1,
\qquad
D'=1000D+\delta,
\]
entonces:
\[
\boxed{
A'=729000A+u1000^{K+1}-\delta\,729\,3^L.
}
\]
Y:
\[
\boxed{
B'=A'+1000^{K+\Delta K}.
}
\]
Esta es la ley exacta de arrastre del cilindro:
\[
\boxed{
\text{el tail/carry se transporta aritméticamente.}
}
\]

\section{Calendario mecánico}

El número de triadas visibles viene dado por:
\[
K(t)=\left\lfloor6(t+1)\log_{1000}3\right\rfloor.
\]
Por tanto:
\[
\Delta K_t=K(t+1)-K(t)\in\{0,1\}.
\]
Los stutters no son anomalías. Son parte del calendario:
\[
\boxed{
\text{abrir }729,\qquad \text{podar }1000.
}
\]

\section{Selector de pares \((u,\delta)\)}

N70 deja muchos candidatos:
\[
(u^\chi,\delta^\chi),
\qquad
\chi\in\{\pi,e,\varphi\}.
\]
La selección real no es por métrica de un solo canal. Se acoplan los tres canales:
\[
B=
\begin{pmatrix}
u^\pi\\
u^e\\
u^\varphi
\end{pmatrix}.
\]
Y se impone firma:
\[
q=\mathbf1^TB,
\]
\[
a=(1,1,-1)B,
\]
\[
c=\frac{\mathbf1^TB-q}{3},
\]
\[
\operatorname{colw},
\qquad
\rho_W=(BA_W)\mathbf1,
\qquad
r=B\mathbf1.
\]
Así:
\[
\boxed{
\text{el selector es acoplado y de frontera, no métrico por canal.}
}
\]

\section{Supervivencia coinductiva}

La firma local no siempre es única. La solución es leer por supervivencia:
\[
S_t=\text{cilindro exacto}+\text{firma de frontera}.
\]
Una rama \(s\in S_t\) sobrevive a profundidad \(h\) si puede extenderse \(h\) fronteras futuras compatibles.

La lectura fuerte:
\[
\boxed{
\text{número HMT}=
\text{rama de límite inverso que sobrevive a todas las fronteras.}
}
\]
Esto explica por qué medir más restringe el futuro:
\[
\boxed{
\text{la precisión es poda de ramas futuras.}
}
\]

\section{Puerta de orientación}

Para una frontera:
\[
q=u^\pi+u^e+u^\varphi,
\qquad
a=u^\pi+u^e-u^\varphi.
\]
Entonces:
\[
q-a=2u^\varphi.
\]
En \(\mathbb F_3\), \(2^{-1}=2\). Por tanto:
\[
\boxed{
u^\varphi=2(q-a).
}
\]
Y:
\[
\boxed{
u^\pi+u^e=q-u^\varphi.
}
\]
Esto prueba:
\[
\boxed{
\text{si }q,a\text{ están fijos, la ambigüedad sólo puede vivir en el reparto de }\pi/e.
}
\]
Así:
\[
\boxed{
\varphi=\text{eje fijo},
\qquad
\pi/e=\text{espejo móvil}.
}
\]

\section{La novena puerta y la orientación}

La novena puerta se estudia en dos lecturas:
\[
t=8
\]
como noveno bloque de seis trits, y:
\[
t=9
\]
como frontera donde:
\[
K(t)=9.
\]
En ambas aparece:
\[
\boxed{
\varphi\text{ fijo},
}
\]
\[
\boxed{
\pi/e\text{ como par móvil},
}
\]
\[
\boxed{
\text{el futuro selecciona la orientación viva}.
}
\]
La novena puerta no es un stutter. Es una puerta de orientación. Hace visible la ley:
\[
\boxed{
\text{la autoescala queda fija y el cierre/propagación se espeja.}
}
\]

\section{Monodromía de nueve puertas}

Definimos:
\[
g(K)=K\bmod9,\qquad 0\equiv9.
\]
Entonces:
\[
\boxed{
g(K+9)=g(K).
}
\]
Pero:
\[
\boxed{
\mathcal S_{K+9}\neq\mathcal S_K\quad\text{en general}.
}
\]
Esto significa:
\[
\boxed{
\text{la torre no es periódica; es monodrómica.}
}
\]
La frase central:
\[
\boxed{
\text{la raíz vuelve; el borde recuerda.}
}
\]
La puerta \(10\) repite la fase de la puerta \(1\):
\[
10\equiv1\pmod9.
\]
Pero no repite la misma frontera. La frontera vuelve cargada por:
\[
\text{carry},\quad
\text{cola},\quad
\text{supervivencia},\quad
\text{orientación}.
\]
Por eso:
\[
\boxed{
\text{la torre no es periódica; es recursiva/fractal.}
}
\]

\section{Fractalización}

La fractalización no es metáfora visual. Cada nueva capa abre:
\[
3^6=729
\]
posibilidades. La lectura decimal poda por:
\[
1000.
\]
Como:
\[
1000\equiv1\pmod9,
\]
cada poda retorna a la unidad de fase, pero el carry conserva memoria.

La forma compacta:
\[
\boxed{
\mathcal E_\infty
=
\text{célula }729
+
\text{poda }1000
+
\text{carry}
+
\text{supervivencia}
+
\text{nueve puertas}.
}
\]
Hacia dentro:
\[
\text{decimales},\quad
\text{cilindros},\quad
\text{límites inversos}.
\]
Hacia fuera:
\[
W_{12},\quad
W_{24},\quad
\Leech,\quad
V^\natural,\quad
\Monster.
\]
Así:
\[
\boxed{
\text{la misma célula se autolee hacia dentro como decimal y hacia fuera como simetría.}
}
\]

\section{Interpretación de la medida}

La comparación clásica con \(\sqrt2\) es instructiva. La diagonal del cuadrado se representa geométricamente, pero no como fracción exacta \(a/b\). Cuanto más se mide, más se precisa su cilindro numérico.

En HMT:
\[
\boxed{
\text{medir}=\text{reducir ramas compatibles de un cilindro.}
}
\]
Por eso la precisión tiene relación con predicción:
\[
\boxed{
\text{más precisión}=\text{menos futuros compatibles}.
}
\]
La medición no sólo describe el presente; poda la frontera de posibilidades.

\section{Pre-continuo y emergencia}

La estructura HMT se sitúa antes de la lectura continua:
\[
\text{pre-continuo}
\longrightarrow
\text{cilindro}
\longrightarrow
\text{frontera}
\longrightarrow
\text{lectura real}.
\]
Ese pre-continuo no debe imaginarse como un espacio pequeño dentro del espacio. No es un microespacio local. Es una estructura de compatibilidades antes de que haya emplazamiento métrico.

La frontera donde cambia la dimensión no es un punto espacial. Es una entrefase:
\[
\boxed{
\text{lugar sin lugar: borde de actualización.}
}
\]
En el cuaderno matemático debe formularse así:
\[
\boxed{
\text{el borde HMT es condición de lectura, no región espacial ordinaria.}
}
\]

\section{Qué queda demostrado y qué queda como programa}

Lo consolidado:
\begin{enumerate}[label=\arabic*.]
\item La APP define una célula de fase suma/producto.
\item La célula triádica activa es \(729=3^6\).
\item La lectura decimal se organiza por \(1000=729+271\).
\item Las semillas \(w_6(\pi),w_6(e),w_6(\varphi)\) abren regiones iniciales.
\item El corrector \(\KR\) permite el cierre canónico de \(\alpha\).
\item \(W_{12}=S(5,6,12)\) y el entorno Golay/Witt rigidifican el borde.
\item \(A_W^2=-I\) da estructura compleja ternaria.
\item \(q_\varphi=w_6(\varphi)A_W^{-1}\) da coordenada Witt de autoescala.
\item La poda decimal correcta es intervalar/cilíndrica.
\item La transición de cilindros \(A,B\) es exacta.
\item La selección no es por canal aislado, sino por firma acoplada de frontera.
\item La supervivencia coinductiva resuelve ambigüedades finitas.
\item La ambigüedad local es orientación \(\pi/e\) con \(\varphi\) fijo.
\item La torre funciona como monodromía de nueve puertas.
\end{enumerate}

Lo que debe formularse con disciplina:
\[
\boxed{
\text{la generación infinita completa requiere un teorema coinductivo global.}
}
\]
Pero ya sabemos la forma correcta:
\[
\boxed{
\text{una rama es HMT si y sólo si sobrevive indefinidamente bajo la máquina de nueve puertas.}
}
\]

\section{Definición propuesta de número HMT}

Definición provisional:
\[
\boxed{
\text{un número HMT es una rama de límite inverso en la célula APP--729 que sobrevive a todas las fronteras de la máquina de nueve puertas.}
}
\]
Más explícitamente:
\[
X=(C_0,C_1,C_2,\ldots)
\]
donde cada \(C_t\) es un cilindro compatible con:
\[
\APP,\quad
729,\quad
1000,\quad
\text{carry},\quad
\text{firma},\quad
\text{Witt},\quad
\text{supervivencia}.
\]
La lectura real es:
\[
\bigcap_{t\ge0} C_t.
\]
Si la intersección es un único real, aparece la cinta decimal. Pero la cinta decimal es la sombra, no el objeto primario.

\section{Frases de cierre}

\[
\boxed{
\text{HMT no genera cifras: genera fronteras que se leen como cifras.}
}
\]
\[
\boxed{
\text{la precisión no añade dígitos: poda futuro.}
}
\]
\[
\boxed{
\text{la raíz vuelve; el borde recuerda.}
}
\]
\[
\boxed{
\text{la novena puerta no es magia: es orientación.}
}
\]
\[
\boxed{
\varphi\text{ fija la escala; }\pi/e\text{ se espejan; el futuro selecciona.}
}
\]
\[
\boxed{
\alpha\text{ es cierre real fino de la contracción }(\pi,e,\varphi).
}
\]
\[
\boxed{
\text{el continuo es lectura estabilizada de una estructura finita de borde.}
}
\]

\end{document}
